\section{VẬN DỤNG BÀI HỌC LỊCH SỬ NHẰM HOÀN THIỆN HỆ THỐNG GIÁO DỤC NGHỀ NGHIỆP VIỆT NAM ĐẾN NĂM 2035}

Lịch sử chỉ thực sự phát huy trọn vẹn giá trị khi những bài học của quá khứ được chuyển hóa thành các quyết sách hành động cho hiện tại và tương lai. Bước vào giai đoạn phát triển mới, trước những biến động sâu sắc của bối cảnh toàn cầu và trong nước, việc vận dụng các bài học lịch sử đã được đúc kết có ý nghĩa sống còn nhằm định hình một hệ thống giáo dục nghề nghiệp Việt Nam hiện đại, mở, linh hoạt và hội nhập quốc tế.

\subsection{Bối cảnh mới và những yêu cầu đặt ra cho giáo dục nghề nghiệp}

\subsubsection{Tác động của Cách mạng công nghiệp lần thứ tư, trí tuệ nhân tạo và kinh tế xanh}

Cuộc Cách mạng công nghiệp lần thứ tư với trọng tâm là trí tuệ nhân tạo, điện toán đám mây, dữ liệu lớn, in ba chiều và công nghệ robot tự hành đang làm thay đổi căn bản cấu trúc của thị trường lao động toàn cầu. Hàng triệu công việc truyền thống mang tính lặp đi lặp lại đang bị máy móc và các thuật toán thông minh thay thế với tốc độ chóng mặt. Ngược lại, nền kinh tế đang chứng kiến sự bùng nổ của các vị trí việc làm mới đòi hỏi các kỹ năng lai ghép (kết hợp giữa kỹ năng chuyên môn sâu với kỹ năng số và tư duy sáng tạo), tiêu biểu như kỹ thuật viên vận hành sản xuất vi mạch bán dẫn, kỹ sư thực hành trí tuệ nhân tạo, kỹ thuật viên bảo trì hệ thống xe điện và kỹ thuật viên năng lượng tái tạo.

Bên cạnh làn sóng công nghệ số, cam kết mạnh mẽ của Chính phủ Việt Nam tại Hội nghị Thượng đỉnh về Biến đổi khí hậu của Liên Hợp Quốc (COP26) về việc đạt mức phát thải ròng bằng không vào năm 2050 đang thúc đẩy quá trình chuyển đổi xanh trên toàn bộ nền kinh tế. Sự xuất hiện của các ngành công nghiệp xanh đòi hỏi hệ thống giáo dục nghề nghiệp phải khẩn trương chuyển đổi nội dung chương trình đào tạo, tích hợp tiêu chuẩn kỹ năng xanh vào mọi ngành nghề, từ xây dựng xanh, nông nghiệp công nghệ cao đến logistics tuần hoàn.

\subsubsection{Thách thức từ xu hướng già hóa dân số và bẫy thu nhập trung bình}

Việt Nam đang trải qua thời kỳ cơ cấu dân số vàng nhưng đồng thời cũng là một trong những quốc gia có tốc độ già hóa dân số nhanh nhất thế giới. Theo dự báo của Tổng cục Thống kê, Việt Nam sẽ chính thức bước vào thời kỳ dân số già từ năm 2036. Lợi thế về nguồn lao động trẻ, dồi dào, giá rẻ đang nhanh chóng thu hẹp lại. 

Nếu không kịp thời nâng cao năng suất lao động xã hội thông qua việc trang bị kỹ năng tay nghề trình độ cao cho lực lượng lao động, Việt Nam sẽ phải đối mặt với nguy cơ chưa giàu đã già và rơi vào bẫy thu nhập trung bình. 

Thực tế này đòi hỏi hệ thống giáo dục nghề nghiệp không thể chỉ tập trung vào đối tượng học sinh tốt nghiệp phổ thông truyền thống, mà phải chuyển hướng mạnh mẽ sang phục vụ toàn thể lực lượng lao động xã hội thông qua các khóa đào tạo lại và đào tạo nâng cao thường xuyên, biến triết lý học tập suốt đời thành hiện thực sống động.

\subsection{Mục tiêu và tầm nhìn phát triển theo các văn kiện chiến lược quốc gia}

Định hướng phát triển của giáo dục nghề nghiệp trong giai đoạn tới đã được Đảng và Nhà nước xác định rõ ràng thông qua hai văn kiện kim chỉ nam có tầm vóc lịch sử:

Thứ nhất là \textbf{Chỉ thị số 21-CT/TW} ngày 04 tháng 5 năm 2023 của Ban Bí thư Trung ương Đảng về tiếp tục đổi mới, phát triển và nâng cao chất lượng giáo dục nghề nghiệp đến năm 2030, tầm nhìn đến năm 2045. Chỉ thị đã nâng tầm phát triển giáo dục nghề nghiệp thành một trong những nhiệm vụ chính trị quan trọng hàng đầu; xác định mục tiêu đến năm 2030 thu hút từ năm mươi đến năm mươi lăm phần trăm học sinh sau khi tốt nghiệp trung học cơ sở và trung học phổ thông vào học tại các cơ sở giáo dục nghề nghiệp; đồng thời các cơ sở đào tạo nghề của Việt Nam phải tiếp cận trình độ các nước phát triển trong nhóm ASEAN-4 và G20.

Thứ hai là \textbf{Quyết định số 2239/QĐ-TTg} ngày 30 tháng 12 năm 2021 của Thủ tướng Chính phủ phê duyệt \textit{Chiến lược phát triển giáo dục nghề nghiệp giai đoạn 2021 -- 2030, tầm nhìn đến năm 2045}. Chiến lược đặt ra chỉ tiêu phấn đấu đến năm 2030 tỷ lệ lao động qua đào tạo có bằng cấp, chứng chỉ đạt từ ba mươi lăm đến bốn mươi phần trăm; hình thành khoảng một trăm trường cao đẳng chất lượng cao, trong đó có khoảng bốn mươi trường tiếp cận trình độ các nước tiên tiến trong khu vực và thế giới; đưa chỉ số chất lượng đào tạo nghề nghiệp của Việt Nam thuộc nhóm ba mươi nước dẫn đầu thế giới.

\subsection{Các khuyến nghị chính sách cụ thể kế thừa từ bài học lịch sử}

Để hiện thực hóa các mục tiêu chiến lược trên cơ sở kế thừa sâu sắc những bài học kinh nghiệm lịch sử, báo cáo đề xuất bốn nhóm khuyến nghị chính sách trọng tâm:

\subsubsection{Hoàn thiện thể chế pháp luật và kiện toàn bộ máy quản lý nhà nước tinh gọn, hiệu lực}

Bài học lịch sử về những tổn thất do phân tán thể chế đòi hỏi công tác hoàn thiện khung khổ pháp lý phải được ưu tiên hàng đầu. Cần khẩn trương tổng kết toàn diện mười năm thi hành Luật Giáo dục nghề nghiệp năm 2014 để xây dựng dự thảo Luật Giáo dục nghề nghiệp (sửa đổi) trình Quốc hội. 

Luật mới cần tập trung tháo gỡ các vướng mắc cốt lõi: tạo cơ chế liên thông thực chất giữa giáo dục nghề nghiệp và giáo dục đại học, cho phép công nhận tín chỉ học tập tích lũy từ kinh nghiệm thực tiễn (Recognition of Prior Learning -- RPL); trao quyền tự chủ học thuật, tổ chức và nhân sự toàn diện cho các trường cao đẳng chất lượng cao; và xác lập địa vị pháp lý bình đẳng của văn bằng cao đẳng kỹ thuật với bằng cử nhân trong hệ thống tuyển dụng công.

Đồng thời, đối với bộ máy quản lý nhà nước, việc sắp xếp, tinh gọn các cơ quan quản lý cần kiên quyết quán triệt nguyên tắc một việc chỉ giao cho một cơ quan chủ trì chịu trách nhiệm chính. Dù cơ quan chuyên trách đặt tại Bộ Giáo dục và Đào tạo hay Bộ Lao động -- Thương binh và Xã hội, điều kiện tiên quyết là phải bảo tồn và phát huy được tính đặc thù của đào tạo nghề nghiệp (gắn liền với thị trường việc làm và tiêu chuẩn kỹ năng nghề của doanh nghiệp), tránh nguy cơ "hàn lâm hóa" trường nghề làm mất đi bản sắc thực hành.

\subsubsection{Đổi mới cơ chế tài chính, đẩy mạnh tự chủ và ban hành chính sách ưu đãi đột phá cho doanh nghiệp}

Tài chính là đòn bẩy quyết định chất lượng giáo dục nghề nghiệp. Nhà nước cần chuyển đổi phương thức phân bổ ngân sách từ hình thức cấp phát cào bằng theo định mức sang cơ chế đấu thầu, đặt hàng đào tạo dựa trên kết quả đầu ra và tỷ lệ người học có việc làm đúng ngành nghề. 

Đối với các cơ sở đào tạo công lập, cần ban hành nghị định chuyên biệt về tự chủ tài chính cho trường nghề, cho phép các trường được liên doanh, liên kết, cho thuê tài sản công để tái đầu tư vào xưởng thực hành và trả thu nhập tăng thêm thỏa đáng nhằm giữ chân các nghệ nhân, giáo viên dạy nghề giỏi.

Đối với khu vực doanh nghiệp, cần thể chế hóa cơ chế đối tác ba bên (Nhà nước -- Nhà trường -- Doanh nghiệp) thông qua các chính sách kinh tế có sức nặng:

Cho phép doanh nghiệp được trừ toàn bộ chi phí đầu tư cho cơ sở giáo dục nghề nghiệp (như tặng máy móc, cử chuyên gia sang giảng dạy, cấp học bổng) vào chi phí hợp lý khi tính thuế thu nhập doanh nghiệp.

Giảm thuế giá trị gia tăng và thuế thu nhập doanh nghiệp đối với các doanh nghiệp tiếp nhận sinh viên thực tập từ sáu tháng trở lên và cam kết trả lương thực tập.

Thành lập các Hội đồng kỹ năng ngành (Sector Skills Councils) với sự tham gia chiếm đa số của đại diện giới chủ sử dụng lao động để trực tiếp tham gia xây dựng chuẩn đầu ra và đánh giá công nhận kỹ năng nghề quốc gia.

\subsubsection{Quy hoạch mạng lưới cơ sở đào tạo theo hướng mở, linh hoạt và đẩy mạnh chuyển đổi số}

Bài học khủng hoảng đầu thời kỳ Đổi mới nhắc nhở rằng mạng lưới trường lớp không thể duy trì theo quy mô hành chính cồng kềnh mà phải được quy hoạch lại theo vùng kinh tế trọng điểm và cụm ngành công nghiệp. Cần kiên quyết sáp nhập, giải thể các trung tâm giáo dục nghề nghiệp, trường trung cấp hoạt động kém hiệu quả, không tuyển sinh được; đồng thời dồn nguồn lực ngân sách để xây dựng các trung tâm quốc gia đào tạo và thực hành nghề chất lượng cao tại ba miền Bắc, Trung, Nam.

Song song với việc quy hoạch mạng lưới vật lý, cần đẩy mạnh chuyển đổi số toàn diện hoạt động đào tạo nghề:

Xây dựng kho học liệu số dùng chung quốc gia gồm các bài giảng thực hành, mô hình mô phỏng ba chiều và ngân hàng câu hỏi đánh giá kỹ năng nghề.

Đầu tư thí điểm các phòng xưởng thực hành ảo ứng dụng công nghệ thực tế ảo (VR) và thực tế ảo tăng cường (AR) để hỗ trợ đào tạo các ngành nghề kỹ thuật cao, độc hại hoặc có chi phí thực hành tiêu hao quá đắt đỏ (như kỹ thuật hàn ngầm dưới nước, điều khiển robot công nghiệp, chẩn đoán lỗi động cơ máy bay).

Ứng dụng trí tuệ nhân tạo để cá nhân hóa tiến trình học tập của sinh viên và hỗ trợ giáo viên theo dõi, đánh giá độ chính xác trong từng thao tác kỹ thuật của người học.

\subsubsection{Nâng cao vị thế xã hội của giáo dục nghề nghiệp qua hệ thống hướng nghiệp sớm và tôn vinh lao động giỏi}

Để giải quyết tận gốc rễ rào cản định kiến bằng cấp đã tồn tại hàng thế kỷ, Nhà nước cần triển khai một chiến lược truyền thông xã hội dài hạn và sâu rộng:

Cần đưa giáo dục hướng nghiệp và trải nghiệm nghề nghiệp thành môn học bắt buộc, thực chất ngay từ bậc trung học cơ sở theo Đề án giáo dục hướng nghiệp và định hướng phân luồng học sinh phổ thông (Quyết định số 522/QĐ-TTg của Thủ tướng Chính phủ). Cần tạo điều kiện cho học sinh phổ thông được trực tiếp tham quan các nhà xưởng hiện đại, tiếp xúc với các cựu sinh viên trường nghề thành đạt để khơi dậy niềm đam mê lao động kỹ thuật từ sớm.

Cải cách căn bản chính sách tiền lương và vị trí việc làm trong khu vực công và khu vực doanh nghiệp nhà nước: thực hiện việc trả lương, bổ nhiệm cán bộ căn cứ vào năng lực thực thi và tiêu chuẩn kỹ năng nghề thực tế chứ không căn cứ vào tên gọi của văn bằng đại học hay cao đẳng.

Định kỳ tổ chức trọng thể các phong trào tôn vinh người thợ giỏi, "Bàn tay vàng", Ngày Kỹ năng lao động Việt Nam (ngày 04 tháng 10 hàng năm) để nâng cao lòng tự hào nghề nghiệp, khẳng định vị thế cao quý của những người trực tiếp làm ra của cải vật chất cho xã hội.
